\section{A sharp extremum for two-dimensional graph towers}
\label{sec:moment-extremum}

The earlier incoming report
\emph{Cyclotomic valuations in graph towers and progression-free sets over finite rings}
proved a finite-field moment formula for spanning-tree valuations and explicitly
asked for the largest value of its coefficient when the quadratic moment
vanishes identically.\cite{priorGraph} We resolve that extremal problem for every even number
of steps in an infinite congruence family of primes. The result is a new
deduction relative to the inspected reports; no claim of historical priority
is made. It uses no RH or GRH assumption.

Let $p>2m$ be prime, and let $s_1,\dots,s_m$ be nonzero vectors of
$\mathbb F_p^2$ spanning that space, with repetitions allowed. Put
\[
 M_j(X,Y)=\sum_{i=1}^m(Xs_{i,1}+Ys_{i,2})^{2j},\qquad
 r(L)=\min\{j\ge1:M_j(L)\ne0\},\qquad
 A(S)=2\sum_{L\in\mathbb P^1(\mathbb F_p)}r(L).
\]
Spanning implies $m\ge2$. The zero/nonzero condition is independent of the representative of $L$.
Newton's identities show that $r(L)\le m$: otherwise all elementary
symmetric functions of the $m$ numbers $(L\cdot s_i)^2$ would vanish,
forcing every projection to be zero, contrary to spanning.

\begin{theorem}[Universal moment bound and an exact even-step maximum]
\label{thm:moment-maximum}
Write $k=\lfloor m/2\rfloor$. Then
\[
                 A(S)\le 2k\bigl(p+1+2m-2k\bigr).
\]
If $m=2k$ is even and $p\equiv1\pmod{2m}$, then
\[
             \boxed{\ \max_S A(S)=m(p+1+m).\ }
\]
For $m\ge4$, the same maximum is attained within the subclass $M_1=0$
identically. Thus this also answers the earlier report's degenerate-quadratic
extremal question in these parameters.
\end{theorem}

The pointwise bound also guarantees the existence of a nonzero moment polynomial of index at most $m$. The main point is that a two-dimensional spanning list cannot make its first
too many even moment \emph{polynomials} vanish simultaneously.

\begin{lemma}[First nonzero moment polynomial]
\label{lem:moment-first}
If $r=\min\{j\ge1:M_j\ne0\text{ as a polynomial}\}$, then
$r\le\lfloor m/2\rfloor$.
\end{lemma}
\begin{proof}
We first record a standard Vandermonde fact. For $d$ distinct projective
linear forms $\ell_1,\dots,\ell_d$, their $N$th powers are linearly independent
if $d\le N+1$ and $N<p$. Choose coordinates making all their slopes finite;
this is possible because $d\le m<p+1$. After rescaling write the forms as
$X+t_aY$. The first $d$ coefficient equations have matrix
$\binom Nj t_a^j$ ($0\le j<d$), a Vandermonde matrix times invertible row
factors. This proves the fact.

Suppose, for a contradiction, that $m<2r$. Partition the $m$ linear forms
$Xs_{i,1}+Ys_{i,2}$ into their $d\ge2$ projective directions. Within direction
$a$, write them as $\lambda_{a,i}\ell_a$, with every $\lambda_{a,i}\ne0$.
Since $M_{r-1}=0$ and $2r-2\ge m-1\ge d-1$, independence gives
\[
             \sum_i\lambda_{a,i}^{2(r-1)}=0
                         \quad\text{for every }a.
\]
Here $r\le m$ by the pointwise Newton argument, so $2r-2<p$ as needed.
Each group therefore has at least two entries, and $2\le d\le m/2<r$.
Set $J=\lfloor d/2\rfloor$. For every $J\le j\le r-1$, the powers
$\ell_a^{2j}$ are independent, because $d\le2j+1$ and $2j<p$. Therefore
\[
                 \sum_i\lambda_{a,i}^{2j}=0
                     \qquad(J\le j\le r-1).
\]
Group the distinct nonzero values $\lambda_{a,i}^2$ in one direction as
$\alpha_1,\dots,\alpha_t$, with positive integer multiplicities
$c_1,\dots,c_t\le m<p$. If $t\le r-J$, the first $t$ equations just displayed
form an invertible Vandermonde system
$\sum_\nu c_\nu\alpha_\nu^{J+j}=0$ ($0\le j<t$), forcing all $c_\nu=0$
in $\mathbb F_p$. This contradicts their range. Every direction thus has
at least $r-J+1$ entries, giving
\[
               m\ge d\bigl(r-\lfloor d/2\rfloor+1\bigr)\ge2r.
\]
For completeness, the last inequality is equality at $d=2$. For $d\ge3$,
use $r\ge d+1$ and $\lfloor d/2\rfloor\le d/2$ to bound its difference
from $2r$ below by $d^2/2-2>0$. This contradicts $m<2r$.
\end{proof}

\begin{proof}[Proof of Theorem~\ref{thm:moment-maximum}]
Let $r$ be as in Lemma~\ref{lem:moment-first}. The nonzero binary homogeneous
polynomial $M_r$ has at most $2r$ projective zeros. At its other projective
points $r(L)=r$, and at each zero $r(L)\le m$. Consequently
\[
 A(S)\le 2r(p+1)+2r(2m-2r)
         =2r(p+1+2m-2r).
\]
This expression increases strictly with integer $r$ for $1\le r\le k$:
the successive difference is $2(p+1+2m-4r-2)>0$ in that range because
$p>2m$. The lemma gives the asserted bound.

For attainment let $m=2k$ and choose $g\in\mathbb F_p^\times$ of order
$2k=m$. Take the $2k$ steps
\[
 (g^i,0),\quad(0,g^i),\qquad 0\le i<k.
\]
Then $g^2$ has order $k$, so
\[
 M_j(X,Y)=
 \begin{cases}
  0,&k\nmid j,\\
  k(X^{2j}+Y^{2j}),&k\mid j.
 \end{cases}
\]
Because $2m=4k$ divides $p-1$, the equation $X^{2k}+Y^{2k}=0$
has exactly $2k=m$ projective roots, all with both coordinates nonzero.
At each such root the next nonzero moment is
$M_{2k}=2kX^{4k}\ne0$, so $r(L)=2k=m$. At the other $p+1-m$ points,
$r(L)=k=m/2$. Thus
\[
  A(S)=2k(p+1-2k)+4k(2k)=m(p+1+m).
\]
For $k\ge2$ the formula also gives $M_1=0$ identically.
\end{proof}

\subsection{Consequences for the actual spanning-tree count}
Lift the steps to integer vectors, and let $\Gamma_n(S)$ be the Cayley
multigraph on $(\mathbb Z/p^n\mathbb Z)^2$ with an undirected edge from
$v$ to $v+s_i$ for each $v$ and each $i$. Keep parallel edges. Spanning
modulo $p$ guarantees connectedness at every level. Write $\kappa_n(S)$
for its number of spanning trees and normalize $v_p(p)=1$.

The earlier moment formula, whose short proof is recalled here, gives
\[
       v_p(\kappa_n(S))=A(S)\frac{p^n-1}{p-1}-2n.
\]
For a primitive character of conductor $p^h$, with $\zeta$ a primitive
$p^h$th root, the corresponding Laplacian eigenvalue is
$\sum_i(2-\zeta^{u\cdot s_i}-\zeta^{-u\cdot s_i})$.
Set $T=\zeta-1$, so $v_p(T)=1/\varphi(p^h)$. Modulo $p$, the invertible
truncated substitution $1+T=\exp Z$ shows that the first nonzero term has
degree $2r(L)\le2m\le p-1$. Its constant and linear coefficients vanish
integrally. Thus its valuation is exactly $2r(L)/\varphi(p^h)$, including
the possible boundary $2r(L)=p-1$. There are $p^{h-1}$ projective lifts
of each residual direction, each orbit having $\varphi(p^h)$ conjugate
characters. Summing their valuations and subtracting $2n$ by the
matrix-tree theorem gives the displayed formula.

\begin{corollary}[Exact maximum at every layer]
For even $m$ and primes $p>2m$ with $p\equiv1\pmod{2m}$,
\[
 \boxed{\quad
 \max_S v_p(\kappa_n(S))
       =m(p+1+m)\frac{p^n-1}{p-1}-2n
       \qquad(n\ge1).
 \quad}
\]
The same fixed root-of-unity step construction attains the maximum
simultaneously at every level, independently of the chosen integral lifts.
\end{corollary}

For example, $m=6,p=13$ gives $A_{\max}=120$ and
\[
               \max_S v_{13}(\kappa_n(S))=10(13^n-1)-2n.
\]
An attaining step list is $(1,0),(4,0),(3,0),(0,1),(0,4),(0,3)$.
Its values at the first three levels are $118,1676,21954$.
At $m=4,p=17$ the maximum coefficient is $88$ and the first-level
valuation is $86$.

The general graph-tower growth theory and the use of finite-field moment
coefficients have substantial precedents.\cite{graphDV,graphDLRV}

\paragraph{Scope.}
This resolves an explicitly posed extremal subproblem in the earlier
ProveIt report and sharpens a family of graph-Iwasawa computations. It does
not prove a number-field Iwasawa conjecture or a nonabelian RH statement.
For odd $m$, the upper bound above remains valid; sharpness is not established
here. Nor is the exact maximum determined here for primes outside
$p\equiv1\pmod{2m}$. The imported finite-field moment identity was already
proved in the earlier report; the new step is its extremal analysis.
